
\documentclass{article}
\usepackage{colortbl}
\usepackage{makecell}
\usepackage{multirow}
\usepackage{supertabular}

\begin{document}

\newcounter{utterance}

\centering \large Interaction Transcript for game `wordle', experiment `high\_frequency\_words\_no\_clue\_no\_critic', episode 4 with reflexion{-}Qwen3.5{-}27B.
\vspace{24pt}

{ \footnotesize  \setcounter{utterance}{1}
\setlength{\tabcolsep}{0pt}
\begin{supertabular}{c@{$\;$}|p{.15\linewidth}@{}p{.15\linewidth}p{.15\linewidth}p{.15\linewidth}p{.15\linewidth}p{.15\linewidth}}
    \# & \multicolumn{2}{c}{Player} && \multicolumn{2}{c}{Game Master} \\
    \hline

    \theutterance \stepcounter{utterance}  
    & & & \multicolumn{4}{p{0.6\linewidth}}{
        \cellcolor[rgb]{0.9,0.9,0.9}{
            \makecell[{{p{\linewidth}}}]{
                \texttt{\tiny{[P1$\langle$GM]}}
                \texttt{You have attempted this same task before and did not succeed. These are the plans you wrote after those attempts.} \\
\texttt{Trial \#1: New plan: The previous attempt failed because I entered a repetitive loop in the explanation text, repeatedly stating that "hence" was invalid due to the letter 'c' being eliminated, instead of immediately selecting a valid alternative word. This violated the rule to provide a concise explanation and a valid guess. Next time, when I identify that a candidate word contains an eliminated letter, I will immediately discard it and select a new word that strictly adheres to the constraints (e.g., 'r' in position 1, 3, or 5; 'e' in position 4; no 's', 'l', 'a', 't', 'c', 'o', 'y', 'd', 'i'). For instance, I would choose "brine" or "grime" (checking 'g' and 'm' are not eliminated) or "prone" (checking 'p' and 'n') and output the explanation and guess in a single, clean pass without self{-}correction loops.} \\
\\ 
\\ 
\texttt{You are a language wizard who likes to guess words by using the given rules.} \\
\\ 
\texttt{Welcome to Wordle! You have six attempts to guess the target word, a valid English word of five lowercase letters (a{-}z). Please use the tags "explanation:" and "guess:" to provide a concise explanation for each guess.} \\
\\ 
\texttt{For instance, if your guess is "apple", your response should be} \\
\texttt{explanation: this is a common five{-}letter English word, and I am starting my guess with this word.} \\
\texttt{guess: apple} \\
\\ 
\texttt{After each guess, your answer will be validated, and you will receive feedback indicating which letters are correct (green), which letters are correct but in the wrong position (yellow), and which letters are incorrect (red). This feedback can be useful in determining which letters to include or exclude in your next guess.} \\
\\ 
\texttt{For example, the feedback for "apple" might be:} \\
\texttt{guess\_feedback: a<yellow> p<yellow> p<green> l<yellow> e<red>} \\
\\ 
\texttt{The explanation should contain details about how the guess\_feedback is used to arrive at a new guess.} \\
\\ 
\texttt{Do not generate any other text.} \\
\\ 
\texttt{Let's begin with your first guess.} \\
            }
        }
    }
     \\ \\

    \theutterance \stepcounter{utterance}  
    & \multicolumn{4}{p{0.6\linewidth}}{
        \cellcolor[rgb]{0.561,0.737,0.561}{
            \makecell[{{p{\linewidth}}}]{
                \texttt{\tiny{[P1$\rangle$GM]}}
                \texttt{explanation: I am starting with a common five{-}letter word containing frequent vowels and consonants to gather initial feedback on letter presence and positions.} \\
\texttt{guess: slate} \\
            }
        }
    }
    & & \\ \\

    \theutterance \stepcounter{utterance}  
    & & & \multicolumn{2}{p{0.3\linewidth}}{
        \cellcolor[rgb]{0.9,0.9,0.9}{
            \makecell[{{p{\linewidth}}}]{
                \texttt{\tiny{[GM$|$GM]}}
                \texttt{attempts = 1} \\
\texttt{target\_word = never} \\
\texttt{guess = slate} \\
\texttt{guess\_feedback = s<red> l<red> a<red> t<red> e<yellow>} \\
            }
        }
    }
    & & \\ \\

    \theutterance \stepcounter{utterance}  
    & & & \multicolumn{4}{p{0.6\linewidth}}{
        \cellcolor[rgb]{0.9,0.9,0.9}{
            \makecell[{{p{\linewidth}}}]{
                \texttt{\tiny{[P1$\langle$GM]}}
                \texttt{guess\_feedback: s<red> l<red> a<red> t<red> e<yellow>} \\
\\ 
\texttt{Provide your response only in this format.} \\
\texttt{explanation: details} \\
\texttt{guess: word} \\
            }
        }
    }
     \\ \\

    \theutterance \stepcounter{utterance}  
    & \multicolumn{4}{p{0.6\linewidth}}{
        \cellcolor[rgb]{0.561,0.737,0.561}{
            \makecell[{{p{\linewidth}}}]{
                \texttt{\tiny{[P1$\rangle$GM]}}
                \texttt{explanation: The feedback indicates 's', 'l', 'a', and 't' are not in the word, while 'e' is present but not in the fifth position. I will choose a word with 'e' in a different position (like the second or fourth) and use common remaining letters like 'r', 'i', 'n', and 'o' to test new possibilities.} \\
\texttt{guess: rinse} \\
            }
        }
    }
    & & \\ \\

    \theutterance \stepcounter{utterance}  
    & & & \multicolumn{2}{p{0.3\linewidth}}{
        \cellcolor[rgb]{0.9,0.9,0.9}{
            \makecell[{{p{\linewidth}}}]{
                \texttt{\tiny{[GM$|$GM]}}
                \texttt{attempts = 2} \\
\texttt{target\_word = never} \\
\texttt{guess = rinse} \\
\texttt{guess\_feedback = r<yellow> i<red> n<yellow> s<red> e<yellow>} \\
            }
        }
    }
    & & \\ \\

    \theutterance \stepcounter{utterance}  
    & & & \multicolumn{4}{p{0.6\linewidth}}{
        \cellcolor[rgb]{0.9,0.9,0.9}{
            \makecell[{{p{\linewidth}}}]{
                \texttt{\tiny{[P1$\langle$GM]}}
                \texttt{guess\_feedback: r<yellow> i<red> n<yellow> s<red> e<yellow>} \\
\\ 
\texttt{Provide your response only in this format.} \\
\texttt{explanation: details} \\
\texttt{guess: word} \\
            }
        }
    }
     \\ \\

    \theutterance \stepcounter{utterance}  
    & \multicolumn{4}{p{0.6\linewidth}}{
        \cellcolor[rgb]{0.561,0.737,0.561}{
            \makecell[{{p{\linewidth}}}]{
                \texttt{\tiny{[P1$\rangle$GM]}}
                \texttt{explanation: The feedback shows 'r', 'n', and 'e' are in the word but in wrong positions, while 'i' and 's' are eliminated. 'r' is not in position 1, 'n' is not in position 3, and 'e' is not in position 5. I need a word with 'r', 'n', 'e' where 'r' is in 2, 3, or 4; 'n' is in 1, 2, 4, or 5; and 'e' is in 1, 2, 3, or 4. 'grime' fits well: 'r' in 2, 'n' is not in 'grime' (wait, 'grime' has no 'n'). Let's try 'grine' (not a word). How about 'brine'? 'b' is new, 'r' in 2 (valid), 'i' is eliminated (invalid). How about 'crone'? 'c' is new, 'r' in 2, 'o' new, 'n' in 4 (valid), 'e' in 5 (invalid, 'e' was yellow in pos 5). 'e' cannot be in 5. Let's try 'green'? 'g' new, 'r' in 2, 'e' in 3 and 4. 'e' in 3 or 4 is valid. 'n' in 5 is valid. 'green' has 'e' twice. Is 'e' in the target twice? We only know one 'e' is present. Let's try 'prone'. 'p' new, 'r' in 2, 'o' new, 'n' in 4, 'e' in 5 (invalid). 'e' cannot be 5. 'r' cannot be 1. 'n' cannot be 3. 'e' cannot be 5. Try 'reign'? 'r' in 1 (invalid). Try 'nerve'? 'n' in 1, 'e' in 2, 'r' in 3, 'v' new, 'e' in 5 (invalid). Try 'nerdy'? 'y' new. 'n' in 1, 'e' in 2, 'r' in 3, 'd' new, 'y' new. 'e' in 2 is valid. 'r' in 3 is valid. 'n' in 1 is valid. 'd' and 'y' are new. Let's check 'nerdy'. 'n' (pos 1, valid), 'e' (pos 2, valid), 'r' (pos 3, valid), 'd' (new), 'y' (} \\
            }
        }
    }
    & & \\ \\

    \theutterance \stepcounter{utterance}  
    & & & \multicolumn{2}{p{0.3\linewidth}}{
        \cellcolor[rgb]{0.9,0.9,0.9}{
            \makecell[{{p{\linewidth}}}]{
                \texttt{\tiny{[GM$|$GM]}}
                \texttt{Error: The response should contain the 'guess:' keyword exactly once.} \\
            }
        }
    }
    & & \\ \\

    \theutterance \stepcounter{utterance}  
    & & & \multicolumn{2}{p{0.3\linewidth}}{
        \cellcolor[rgb]{0.9,0.9,0.9}{
            \makecell[{{p{\linewidth}}}]{
                \texttt{\tiny{[GM$|$GM]}}
                \texttt{game\_result = ABORT} \\
            }
        }
    }
    & & \\ \\

    \theutterance \stepcounter{utterance}  
    & & \multicolumn{4}{p{0.6\linewidth}}{
        \cellcolor[rgb]{0.95,0.95,0.95}{
            \makecell[{{p{\linewidth}}}]{
                \texttt{\tiny{[P1$\leftrightarrow$P]}}
                \texttt{Post{-}episode Reflexion (not sent to GM):} \\
\texttt{New plan: The previous attempt failed because I included internal monologue, self{-}correction, and incomplete sentences (e.g., "wait, 'grime' has no 'n'", "Let's try...") within the "explanation" field, and the message was cut off before providing the "guess:" tag. This violated the rule to provide a concise explanation and a valid guess in a specific format. Next time, I will perform all word validation and logic internally before generating the response. I will ensure the explanation is a single, coherent paragraph that directly justifies the final chosen word without showing the trial{-}and{-}error process. I will strictly verify that the chosen word is a valid five{-}letter English word, contains no eliminated letters (s, l, a, t, i), and places the known letters (r, n, e) in valid positions (r not in 1, n not in 3, e not in 5). Finally, I will ensure the response ends with the "guess:" tag followed immediately by the word, with no trailing text or incomplete thoughts. For example, based on the constraints, I might select "nerve" (if 'v' is allowed) or "green" (if double 'e' is possible) or "crone" (checking 'c' and 'o' are new), but I will only output the final decision. A valid candidate fitting r (not 1), n (not 3), e (not 5) with no s,l,a,t,i could be "brent" (t is eliminated), "crepe" (c,r,e,p,e {-} r in 2, e in 3 and 5? e in 5 is bad). Let's try "renew" (r in 1 bad). "nerdy" (d,y new, n in 1, e in 2, r in 3). This fits. I will output: explanation: Based on feedback, 's', 'l', 'a', 't', 'i' are eliminated. 'r', 'n', 'e' are present but misplaced. 'r' is not in pos 1, 'n' not in pos 3, 'e' not in pos 5. I choose 'nerdy' which places 'n' in 1, 'e' in 2, 'r' in 3, and introduces new letters 'd' and 'y'. guess: nerdy} \\
            }
        }
    }
    & \\ \\

\end{supertabular}
}

\end{document}
